\section{Separation on the translated quarter mesh}
\label{sec:area_law_quarter_mesh}

The numerical estimates used in the initial-star argument come from a
common translated lattice. We record them separately from the construction
of the active interfaces; see \cite[Section~11]{OpenAI2026AreaLaw}.

% Source: 10-geometry.tex, geometry:initial-stars, lines 352--359.
% These estimates concern actual marked points and allowed affine lines.
% Local layer exclusion is included below; active fan interfaces and
% isolated stars remain separate.

\begin{theorem}[Nonadjacent layers avoid fine neighborhoods]
    \label{thm:area_law_local_layer_exclusion}
    \lean{TNLean.PEPS.AreaLaw.Geometry.dyadicLayer_dist_nonadjacent_fineScale}
    \leanok
    \uses{def:area_law_dyadic_layers,def:area_law_fine_belt_scales}
    Let $C_0\ge2$ and $k\ge50{,}000{,}000$, and put $t_k=2^{\ell_k}$.
    If $x\in\overline{D_k}$, $y\in\overline{D_h}$, and the layer
    indices differ by at least two, then
    \begin{align}
        \|x-y\|_\infty & \ge16t_k.
    \end{align}
    The threshold is uniform in the finite endpoint set and origin.
\end{theorem}
\begin{proof}\leanok
    \uses{thm:area_law_dyadic_nonadjacent,thm:area_law_scale_side_ratios}
    The scale-ratio estimate with exponent gap five gives $32t_k\le2^k$.
    The nonadjacent-layer bound gives
    $\|x-y\|_\infty\ge(C_0-1)2^{\max(k,h)}/2\ge2^k/2\ge16t_k$.
\end{proof}

\begin{theorem}[Only neighboring layers occur nearby]
    \label{thm:area_law_local_layer_indices}
    \lean{TNLean.PEPS.AreaLaw.Geometry.dyadicLayer_nearby_indices}
    \leanok
    \uses{def:area_law_dyadic_layers,def:area_law_fine_belt_scales}
    Under the same conditions on $C_0$ and $k$, if
    $x\in\overline{D_k}$, $y\in\overline{D_h}$ and
    $\|x-y\|_\infty<10t_k$, then
    \begin{align}
        h & \le k+1, & k & \le h+1.
    \end{align}
\end{theorem}
\begin{proof}\leanok
    \uses{thm:area_law_local_layer_exclusion}
    A difference of at least two would give distance at least $16t_k$,
    contradicting the strict upper bound since $t_k>0$.
\end{proof}

\begin{definition}[Translated planar mesh]
    \label{def:area_law_affine_mesh}
    \lean{TNLean.PEPS.AreaLaw.Geometry.affineMesh}
    \leanok
    For $o\in\mathbb R^2$ and $q\in\mathbb R$, define
    \begin{align}
        \mathcal M(o,q) & =\{o+qz:z\in\mathbb Z^2\}.
    \end{align}
\end{definition}

\begin{theorem}[Actual marks lie on the quarter mesh]
    \label{thm:area_law_belt_marks_quarter_mesh}
    \lean{TNLean.PEPS.AreaLaw.Geometry.beltMarks_subset_affineMesh}
    \leanok
    \uses{def:area_law_affine_mesh,def:area_law_belt_marks}
    Let $j,\ell$ be nonnegative integers with $\ell\le j+1$.
    For every finite cell-index set $F$,
    \begin{align}
        M_j(F) & \subseteq\mathcal M(o,2^\ell/4).
    \end{align}
    In particular, cells of sides $t/2$, $t$ and $2t$, where
    $t=2^\ell$ and $\ell\ge1$, have all their marks on the same
    mesh $\mathcal M(o,t/4)$.
\end{theorem}
\begin{proof}\leanok
    A mark of a cell of side $2^j$ is $o+2^j(z+e/2)$, where
    $e\in\{0,1,2\}^2$. Since $\ell\le j+1$, the vector
    $2^{j+1-\ell}(2z+e)$ is integral, and the mark equals
    $o+(2^\ell/4)2^{j+1-\ell}(2z+e)$.
\end{proof}

\begin{theorem}[Distinct mesh points are separated]
    \label{thm:area_law_affine_mesh_point_separation}
    \lean{TNLean.PEPS.AreaLaw.Geometry.affineMesh_dist_ge}
    \leanok
    \uses{def:area_law_affine_mesh}
    If $q>0$ and $x,y\in\mathcal M(o,q)$ are distinct, then
    \begin{align}
        \|x-y\|_\infty & \ge q.
    \end{align}
\end{theorem}
\begin{proof}\leanok
    At least one coordinate of the integer vectors indexing $x$ and
    $y$ differs. Its difference has absolute value at least one, so
    the corresponding coordinate difference of $x$ and $y$ has
    absolute value at least $q$.
\end{proof}

\begin{theorem}[Clearance from a nonincident allowed line]
    \label{thm:area_law_affine_mesh_line_clearance}
    \lean{TNLean.PEPS.AreaLaw.Geometry.affineMesh_line_dist_ge}
    \leanok
    \uses{def:area_law_affine_mesh,def:area_law_template}
    Suppose $q>0$ and $u,v,x\in\mathcal M(o,q)$. Let $L$ be the
    affine span of $u,v$. Suppose either $u=v$, or $L$ has slope
    $0$, $\infty$, $1$ or $-1$. If $x\notin L$, then
    \begin{align}
        \|x-y\|_\infty & \ge q/2\qquad(y\in L).
    \end{align}
    Thus on a quarter mesh the clearance is at least $t/8$.
\end{theorem}
\begin{proof}\leanok
    \uses{thm:area_law_affine_mesh_point_separation}
    If $u=v$, use the point-separation bound. Otherwise, the supporting
    line has a constant coordinate, coordinate sum or coordinate
    difference. Its value at a missing mesh point differs from the
    line's value by a nonzero integer multiple of $q$. A coordinate
    difference is bounded by the sup distance; a difference of sums
    or differences is bounded by twice that distance. The result follows.
\end{proof}

\begin{theorem}[Clearance from a nonincident allowed segment]
    \label{thm:area_law_affine_mesh_segment_clearance}
    \lean{TNLean.PEPS.AreaLaw.Geometry.affineMesh_segment_dist_ge_of_not_mem}
    \leanok
    \uses{def:area_law_affine_mesh,def:area_law_template}
    Let $q>0$ and $a,b,p\in\mathcal M(o,q)$. Suppose that the segment
    $[a,b]$ is horizontal, vertical or diagonal, with coincident endpoints
    also permitted. If $p\notin[a,b]$, then every $x\in[a,b]$ satisfies
    $\|p-x\|_\infty\ge q/2$.
    This estimate is used in the isolation arguments for the initial stars
    and recursive repair in \cite[Section~11]{OpenAI2026AreaLaw}.
\end{theorem}
\begin{proof}\leanok
    \uses{thm:area_law_affine_mesh_point_separation,thm:area_law_affine_mesh_line_clearance}
    If $a=b$, then $x=a\ne p$, so mesh-point separation gives the
    stronger bound $q$. Otherwise let $L$ be the supporting line of
    $[a,b]$. If $p\notin L$, apply the line-clearance estimate.
    If $p\in L\setminus[a,b]$, one endpoint $e\in\{a,b\}$ lies
    between $p$ and $x$. Since $p\ne e$, mesh-point separation and
    distance additivity give
    \begin{align}
        q & \le\|p-e\|_\infty
        \le\|p-e\|_\infty+\|e-x\|_\infty
        =\|p-x\|_\infty.
    \end{align}
    In both cases the asserted half-spacing bound follows.
\end{proof}
